% Generated draft: requires mathematical and source review.
\documentclass[11pt]{article}
\usepackage[margin=25mm]{geometry}
\usepackage{fontspec}
\setmainfont{Latin Modern Roman}
\usepackage{amsmath,amssymb,mathtools,unicode-math}
\setmathfont{Latin Modern Math}
\usepackage{enumitem,needspace,xurl,hyperref}
\setlength{\parindent}{0pt}
\setlength{\parskip}{5pt}
\setlength{\emergencystretch}{3em}
\newcommand{\N}{\mathbb{N}}
\newcommand{\Z}{\mathbb{Z}}
\newcommand{\Q}{\mathbb{Q}}
\newcommand{\R}{\mathbb{R}}
\newcommand{\C}{\mathbb{C}}
\newcommand{\F}{\mathbb{F}}
\newcommand{\A}{\mathbb{A}}
\newcommand{\PP}{\mathbb{P}}
\newcommand{\E}{\mathbb{E}}
\newcommand{\eps}{\varepsilon}
\newcommand{\dd}{\,\mathrm{d}}
\begin{document}
\section*{UnsolvedMath ID: 30002315}
\textbf{Source problem code:} catalogue entry
\subsection{Definitions and mathematical statement}
Let \(\mathcal G=(G_n)_{n\geq 1}\) be a sequence of connected graphs with uniformly bounded vertex degrees. The displayed formulas require each \(G_n\) to have finitely many vertices and edges and at least one vertex; these requirements are implicit in the formulas rather than explicitly stated in the supplied record. No assumption that \(|V(G_n)|\to\infty\) is supplied, and none is added here.

For a field \(K\), interpret the cycle space by choosing an orientation of each edge and setting
\[
C_K(G)=\ker\bigl(\partial:K^{E(G)}\longrightarrow K^{V(G)}\bigr),
\qquad \partial(e)=v-u
\]
for an edge oriented from \(u\) to \(v\). For a positive integer \(q\), let \(C_K^q(G)\) be the subspace spanned by the oriented vectors of cycles of length strictly less than \(q\). Define
\[
s_K^q(\mathcal G)=\liminf_{n\to\infty}
\frac{|E(G_n)|-\dim_K C_K^q(G_n)-|V(G_n)|}{|V(G_n)|},
\qquad
\beta_K(\mathcal G)=\inf_{q\in\N,\ q\geq 1}s_K^q(\mathcal G).
\]
Using positive integer cutoffs makes the domain of the supplied infimum explicit. The source does not specify whether loops or multiple edges are allowed, or its precise convention for cycles in those cases.

The supplied record asks the following questions.
\begin{enumerate}
\item Does there exist such a graph sequence \(\mathcal G\) and a prime \(p\) for which the supplied inequality
\[
\beta_{\Q}(\mathcal G)\geq\beta_{\F_p}(\mathcal G)
\]
is strict? The record does not specify whether the question is intended for an arbitrary prescribed prime or for at least one prime. There is also a potential sign inconsistency requiring source review: the usual comparison of ranks of integral cycle matrices over \(\Q\) and \(\F_p\) suggests the opposite weak inequality under the displayed definition. The supplied direction is retained here without certification or correction.
\item Does there exist such a graph sequence \(\mathcal G\) and a field \(K\) for which
\[
\beta_K(\mathcal G)+1<c(\mathcal G),
\]
where \(c(\mathcal G)\) denotes the combinatorial cost intended by the source? The record supplies neither a definition nor a normalization of combinatorial cost. Consequently this question remains conditional on recovering that definition. It is also unclear whether \(K\) is to be prescribed arbitrarily or may be chosen with the example.
\item Let \(\Gamma\) be a finitely generated residually finite group. Does its rank gradient along a Farber chain of normal finite-index subgroups depend on the chain?

To make the terminology explicit, take a descending chain \((\Gamma_n)_{n\geq 1}\) of normal finite-index subgroups. The Farber condition means that for every \(g\in\Gamma\setminus\{1\}\),
\[
\lim_{n\to\infty}
\frac{|\{x\Gamma_n\in\Gamma/\Gamma_n:gx\Gamma_n=x\Gamma_n\}|}
{[\Gamma:\Gamma_n]}=0.
\]
Write \(d(H)\) for the minimum cardinality of a generating set of a finitely generated group \(H\), with \(d(\{1\})=0\). Using the standard normalized convention, set
\[
\operatorname{RG}(\Gamma,(\Gamma_n))
=\inf_{n\geq 1}\frac{d(\Gamma_n)-1}{[\Gamma:\Gamma_n]}.
\]
For a descending chain this infimum is also the limit of the displayed normalized ranks. These definitions spell out standard conventions; the supplied record does not itself give them.

Precisely, is it true that for every finitely generated residually finite \(\Gamma\) and every two normal Farber chains \((\Gamma_n)\) and \((\Lambda_n)\),
\[
\operatorname{RG}(\Gamma,(\Gamma_n))
=\operatorname{RG}(\Gamma,(\Lambda_n))?
\]
Equivalently, the dependence question asks whether there exist such a group and two such chains with different rank gradients.
\end{enumerate}
The record labels the combined problem open, but its supplied literature assessment explicitly discusses only the two graph-invariant strictness questions. No separate literature assessment of the rank-gradient question is supplied.
\subsection{Short English statement}
The record asks whether two bounds involving graph cycle-space invariants can be strict, and whether a finitely generated residually finite group's rank gradient can vary between normal Farber chains. The first bound has a possible sign inconsistency, and the definition of combinatorial cost is missing. These are unresolved source issues. All statements remain drafts requiring human mathematical review.
\subsection{Sources}
\begin{itemize}
\item Supplied local UnsolvedMath registry record, permanent ID \(30002315\), source problem code \texttt{catalogue entry}, in the Oberwolfach Reports Open Problems collection; title: \emph{Strict Graph-Invariant Bounds and Farber-Chain Rank Gradients}. No individual proposer is named. The supplied proposed year is 2013.
\item Supplied report citation: \emph{Arbeitsgemeinschaft: Limits of Structures (2014)}, Oberwolfach Reports, pp.\ 963--1024, DOI \texttt{10.4171/owr/2013/16}. \url{https://doi.org/10.4171/owr/2013/16}. The literature-source label is \emph{Oberwolfach Report 16/2013}. The citation's parenthetical year 2014 differs from its report label, DOI year, and supplied proposed year; this discrepancy is unresolved.
\item Supplied literature source: \emph{Annales Henri Lebesgue}. \url{https://www.numdam.org/item/10.5802/ahl.102.pdf}. No article title, authors, or precise supporting location is supplied.
\item Literature checked at: \texttt{2026-08-21}. This is supplied dated curation, not independently verified. Its supplied status is open, and its assessment says that the checked literature develops the finite beta invariants and related bounds without settling either strictness question in the stated generality. The background explicitly calls this dated literature triage and requests independent verification of time-sensitive status and cited claims.
\item The registry labels the statement \texttt{corrected\_verified} and reports that it was checked against the official Oberwolfach report to restore hypotheses and notation. Those are supplied verification claims; no source visit or independent verification is asserted here.
\end{itemize}
\end{document}
